\documentclass[11pt,a4paper]{article}
\usepackage[T1]{fontenc}
\usepackage[utf8]{inputenc}
\usepackage{lmodern}
\usepackage{amsmath,amssymb}
\usepackage[margin=25mm]{geometry}
\usepackage[hidelinks]{hyperref}
\hypersetup{pdftitle={The bridge midpoint on a compact group at any cut: windings and the centre},pdfauthor={navstokgap research working notes}}
\newcommand{\tightlist}{\setlength{\itemsep}{0pt}\setlength{\parskip}{0pt}}
\setlength{\emergencystretch}{3em}
\setcounter{secnumdepth}{0}
\begin{document}
\hypertarget{the-bridge-midpoint-on-a-compact-group-at-any-cut-windings-and-the-centre}{%
\section{The bridge midpoint on a compact group at any cut: windings and
the
centre}\label{the-bridge-midpoint-on-a-compact-group-at-any-cut-windings-and-the-centre}}

\textbf{Result, 2026-09-28 (Claude; written derivation, to be refereed
by GPT-6 Astra).} Let \(G\) be a compact simply connected Lie group and
take the heat-kernel bridge of total heat time \(t\) from \(e\) to
\(g=e^{iX}\), cut at fraction \(s\in(0,1)\), so that the first piece has
heat time \(st\). For every irreducible representation \(\lambda\) the
midpoint expectation of its character is an exact finite sum, over the
weights \(\kappa\) of \(\lambda\), of image sums over the coroot lattice
\(Q^\vee\) (Theorem 1):

\[E\,\chi_\lambda(m)=\frac{\displaystyle\sum_\kappa m_\lambda(\kappa)\,
e^{-ts(1-s)|\kappa|^2/2}\,e^{is\langle\kappa,X\rangle}
\sum_{H\in Q^\vee}e^{2\pi i(1-s)\langle\kappa,H\rangle}\,
\pi\!\Bigl(\tfrac it(X-2\pi H)-(1-s)\kappa\Bigr)\,G_H}
{\displaystyle\sum_{H\in Q^\vee}\pi\!\Bigl(\tfrac it(X-2\pi H)\Bigr)\,G_H},
\qquad G_H=e^{-|X-2\pi H|^2/(2t)}, \tag{1}\]

with \(m_\lambda(\kappa)\) the weight multiplicities and
\(\pi(v)=\prod_{\alpha>0}\langle v,\alpha\rangle/\langle\rho,\alpha\rangle\)
the Weyl dimension polynomial.

\begin{itemize}
\tightlist
\item
  \textbf{Corollary 2 (reductions).} For \(U(1)\), (1) is the winding
  mixture of \href{series-parallel-gauge-refinement.md}{Corollary
  5\(_s\)}(a). For \(SU(2)\) at \(s=\frac12\) it is Theorems 1 and 3 of
  the \href{su2-midpoint-exact.md}{\(SU(2)\) midpoint note}, image terms
  included.
\item
  \textbf{Corollary 3 (the centre).} The phase
  \(e^{2\pi i(1-s)\langle\kappa,H\rangle}\) of an image is the same for
  all weights of every representation iff \((1-s)H\in P^\vee\), the
  coweight lattice, and it is then the central character of \(\lambda\)
  at \(z_H=e^{2\pi i(1-s)H}\). At \(s=p/q\) in lowest terms these
  central images reach exactly the \(q\)-torsion of the centre:
  \(\mathbb Z_{\gcd(q,N)}\) for \(SU(N)\).
\item
  \textbf{Corollary 4 (\(SU(3)\)).} At \(s=\frac13\) the central images
  are \(H\in3P^\vee\); the smallest are the two Weyl triples of
  \(\pm(1,1,-2)\), of relative weight about \(e^{-24\pi^2/t}\), and they
  carry \(\omega^2\) and \(\omega\) on the fundamental
  (\(\omega=e^{2\pi i/3}\)), the conjugates on the antifundamental and
  \(1\) on the adjoint. At \(s=\frac12\) every image with a nontrivial
  phase has a weight-dependent one: no halving of \(SU(3)\) reaches the
  centre.
\end{itemize}

For \(SU(2)\) at \(s=\frac12\) the phase is \((-1)^w\) for half-integer
spin and \(1\) for integer spin. This is Dirac's belt trick,
\(\pi_1(SO(3))=\mathbb Z_2\)
(\href{https://doi.org/10.1112/jlms/s1-17.3.173}{Newman 1942},
metadata), appearing in the exact midpoint law: a bridge winding \(w\)
times round the group makes \(w\) extra \(720^\circ\) turns, and its
midpoint makes \(w\) extra \(360^\circ\) turns, which the spinor
characters detect. In \(1+3\) this internal \(\mathbb Z_2\) meets the
spatial one in ``spin from isospin''
(\href{https://doi.org/10.1103/PhysRevLett.36.1116}{Jackiw and Rebbi
1976}; \href{https://doi.org/10.1103/PhysRevLett.36.1119}{Hasenfratz and
't Hooft 1976}; metadata). For \(SU(3)\) the same mechanism carries the
triality, and only cuts at thirds (or at denominators divisible by
\(3\)) produce it.

The ingredients are character orthogonality, the Brauer--Klimyk rule and
Poisson summation, as in the \(SU(2)\) note; the denominator of (1) is
the classical image formula for the heat kernel on a compact group
(\href{https://doi.org/10.4310/jdg/1214438176}{Fegan 1983}, metadata).
The same unfolding and coroot Poisson summation give the heat trace at
the identity and its exponential Seeley--DeWitt structure,
\(a_k=(R/6)^k/k!\), in the author's earlier project
(\href{https://github.com/arivero/physres1/blob/9e31f51/paper/main.md}{physres1,
Remark D9.1o\(''\)}, 2026-02, read); at \(X\to0\) the \(H=0\) term of
the denominator of (1) gives
\(k_t(e)\propto t^{-\dim G/2}e^{t|\rho|^2/2}\), the same structure
(\(R/6=|\rho|^2\) in this metric; for \(SU(2)\), radius 2,
\(R/6=\frac14\)). No novelty is claimed for the method.

\hypertarget{conventions}{%
\subsection{1. Conventions}\label{conventions}}

The group metric is that of the
\href{series-parallel-gauge-refinement.md}{series/parallel note} §1:
\(-\Delta_G\chi_\lambda=C_2(\lambda)\chi_\lambda\) with
\(C_2(\lambda)=|\lambda+\rho|^2-|\rho|^2\) in the inner product on
weights dual to the metric on the Cartan algebra \(\mathfrak t\). For
\(SU(N)\) roots have \(|\alpha|^2=1\) and coroots
\(\alpha^\vee=2\alpha/|\alpha|^2\) have \(|\alpha^\vee|^2=4\), which is
the metric \(|Y|^2=2\,{\rm tr}\,Y^2\) on \(\mathfrak t\); then
\(C_2(j)=j(j+1)\) for \(SU(2)\) and \(C_2(\mathbf 3)=\frac43\) for
\(SU(3)\). Torus elements are \(e^{iX}\), \(X\in\mathfrak t\); the
pairing of the weight lattice \(P\) with the coroot lattice \(Q^\vee\)
is integral, \(P\) and \(Q^\vee\) are dual lattices, and \(e^{iX}=e\)
iff \(X\in2\pi Q^\vee\) (\(G\) simply connected). The coweight lattice
\(P^\vee\) is dual to the root lattice \(Q\), and
\(P^\vee/Q^\vee\cong Z(G)\) through \(Y\mapsto e^{2\pi iY}\). For
\(SU(2)\), \(|X|=\theta\) is the rotation angle of the \(SU(2)\) note,
and \(X-2\pi w\alpha^\vee\) has length \(|\theta-4\pi w|\).

Heat kernel
\(k_t=\sum_\lambda d_\lambda e^{-tC_2(\lambda)/2}\chi_\lambda\). The
bridge midpoint \(m\) has density
\(k_{st}(m)\,k_{(1-s)t}(m^{-1}g)/k_t(g)\) with respect to Haar measure;
the bridge law (4) of the series/parallel note becomes this one after a
left translation, and Corollary 2\(_s\) there gives the cut faces heat
times \(st\) and \((1-s)t\).

Weyl: \(A_v(X)=\sum_{w\in W}\varepsilon(w)e^{i\langle wv,X\rangle}\),
\(\chi_\lambda=A_{\lambda+\rho}/A_\rho\),
\(d_\lambda=\pi(\lambda+\rho)\). Extend
\(\chi_\nu:=A_{\nu+\rho}/A_\rho\) to every \(\nu\in P\): it is zero or
\(\pm\) an irreducible character, and
\(C_2(\nu):=|\nu+\rho|^2-|\rho|^2\) is invariant under the shifted Weyl
action, so it is the Casimir of that character.

\emph{Brauer--Klimyk rule.}
\(\chi_\lambda\chi_\mu=\sum_\kappa m_\lambda(\kappa)\chi_{\mu+\kappa}\)
for dominant \(\mu\). Indeed
\(\chi_\lambda A_{\mu+\rho}=\sum_\kappa m_\lambda(\kappa)e^{i\langle\kappa,X\rangle}\sum_w\varepsilon(w)e^{i\langle w(\mu+\rho),X\rangle}\),
and replacing \(\kappa\) by \(w\kappa\), which leaves \(m_\lambda\)
invariant, gives \(\sum_\kappa m_\lambda(\kappa)A_{\mu+\kappa+\rho}\).

\hypertarget{theorem-1-and-its-proof}{%
\subsection{2. Theorem 1 and its proof}\label{theorem-1-and-its-proof}}

\textbf{Theorem 1.} For \(X\) regular, (1) holds; for singular \(X\) it
holds by continuity.

\emph{Proof.} (i) \emph{Characters.} Expanding both heat kernels,
\(k_t(g)E\chi_\lambda(m)=\sum_{\mu,\nu}d_\mu d_\nu e^{-stC_2(\mu)/2-(1-s)tC_2(\nu)/2} \int\chi_\lambda(m)\chi_\mu(m)\chi_\nu(m^{-1}g)\,dm\).
The Brauer--Klimyk rule and
\(\int\chi_a(m)\chi_b(m^{-1}g)dm=\delta_{ab}\chi_a(g)/d_a\) give

\[k_t(g)\,E\chi_\lambda(m)=\sum_{\mu\ {\rm dominant}}d_\mu\sum_\kappa m_\lambda(\kappa)\,
e^{-stC_2(\mu)/2-(1-s)tC_2(\mu+\kappa)/2}\,\chi_{\mu+\kappa}(g).\]

\begin{enumerate}
\def\labelenumi{(\roman{enumi})}
\setcounter{enumi}{1}
\tightlist
\item
  \emph{Unfolding.} Multiply by \(A_\rho(X)\) and put \(v=\mu+\rho\),
  which runs over the strictly dominant weights. The summand becomes
  \(\pi(v)\,m_\lambda(\kappa)\,e^{-\frac t2[s|v|^2+(1-s)|v+\kappa|^2-|\rho|^2]}A_{v+\kappa}(X)\).
  Expand \(A_{v+\kappa}\) and substitute \(v\mapsto w^{-1}v\),
  \(\kappa\mapsto w^{-1}\kappa\): \(\varepsilon(w)\pi(w^{-1}v)=\pi(v)\),
  and \(m_\lambda\) and the norms are Weyl invariant. Each regular
  weight is the image of exactly one strictly dominant weight, and
  \(\pi\) vanishes on the walls, so
\end{enumerate}

\[A_\rho(X)k_t(g)E\chi_\lambda(m)=e^{t|\rho|^2/2}\sum_\kappa m_\lambda(\kappa)
\sum_{v\in P}\pi(v)\,e^{-\frac t2[s|v|^2+(1-s)|v+\kappa|^2]}e^{i\langle v+\kappa,X\rangle}.\]

\begin{enumerate}
\def\labelenumi{(\roman{enumi})}
\setcounter{enumi}{2}
\item
  \emph{Square.}
  \(s|v|^2+(1-s)|v+\kappa|^2=|v+(1-s)\kappa|^2+s(1-s)|\kappa|^2\).
\item
  \emph{Poisson summation} over \(P\), whose dual lattice is \(Q^\vee\):
  \(\sum_{v\in P}f(v)={\rm vol}(\mathfrak t^*/P)^{-1}\sum_{H\in Q^\vee}\hat f(H)\)
  with \(\hat f(H)=\int f(y)e^{-2\pi i\langle y,H\rangle}dy\). With
  \(y=u-(1-s)\kappa\) one has
  \(\langle y+\kappa,X\rangle=\langle u,X\rangle+s\langle\kappa,X\rangle\)
  and
  \(\langle y,H\rangle=\langle u,H\rangle-(1-s)\langle\kappa,H\rangle\),
  so
\end{enumerate}

\[\hat f(H)=e^{is\langle\kappa,X\rangle}e^{2\pi i(1-s)\langle\kappa,H\rangle}\int\pi(u-(1-s)\kappa)\,
e^{-t|u|^2/2}e^{i\langle u,Z\rangle}du,\qquad Z=X-2\pi H.\]

\begin{enumerate}
\def\labelenumi{(\alph{enumi})}
\setcounter{enumi}{21}
\tightlist
\item
  \emph{The Gaussian integral.} \(\pi\) is harmonic: \(\Delta\pi\) is
  Weyl-antisymmetric of lower degree, and every Weyl-antisymmetric
  polynomial is divisible by \(\pi\). Its derivatives are harmonic too.
  For a harmonic homogeneous \(h\) of degree \(d\), Hobson's formula
  \(h(\partial)F(|Z|^2)=2^dh(Z)F^{(d)}(|Z|^2)\) gives
  \(h(-i\partial_Z)e^{-|Z|^2/(2t)}=h(iZ/t)\,e^{-|Z|^2/(2t)}\).
  Taylor-expanding \(\pi(u-c)\) in \(c\) into harmonic pieces and
  summing back,
\end{enumerate}

\[\int\pi(u-c)\,e^{-t|u|^2/2+i\langle u,Z\rangle}du=\Bigl(\frac{2\pi}t\Bigr)^{n/2}
\pi\Bigl(\frac{iZ}t-c\Bigr)e^{-|Z|^2/(2t)},\qquad n=\dim\mathfrak t.\]

\begin{enumerate}
\def\labelenumi{(\roman{enumi})}
\setcounter{enumi}{5}
\tightlist
\item
  The trivial representation (\(\kappa=0\) only) gives the denominator.
  The constants \(e^{t|\rho|^2/2}\), \({\rm vol}(\mathfrak t^*/P)^{-1}\)
  and \((2\pi/t)^{n/2}\) cancel. All sums converge absolutely.
  \(\square\)
\end{enumerate}

\textbf{Theorem 1\('\) (several cuts).} Cut the same bridge at
\(0<s_1<\dots<s_n<1\), with layer values \(m_1,\dots,m_n\), and take
irreducible representations \(\lambda_1,\dots,\lambda_n\). With
\(\boldsymbol\kappa=(\kappa_1,\dots,\kappa_n)\) running over weights of
\(\lambda_1,\dots,\lambda_n\),
\(A(\boldsymbol\kappa)=\sum_j(1-s_j)\kappa_j\) and the bridge covariance
form
\(Q(\boldsymbol\kappa)=\sum_{i,j}\min(s_i,s_j)\,(1-\max(s_i,s_j))\langle\kappa_i,\kappa_j\rangle\),

\[E\prod_j\chi_{\lambda_j}(m_j)=\frac{\displaystyle\sum_{\boldsymbol\kappa}\prod_jm_{\lambda_j}(\kappa_j)\,
e^{-tQ(\boldsymbol\kappa)/2}\,e^{i\langle\sum_js_j\kappa_j,X\rangle}\sum_{H\in Q^\vee}e^{2\pi i\langle A(\boldsymbol\kappa),H\rangle}\,
\pi\!\Bigl(\tfrac it(X-2\pi H)-A(\boldsymbol\kappa)\Bigr)G_H}
{\displaystyle\sum_{H\in Q^\vee}\pi\!\Bigl(\tfrac it(X-2\pi H)\Bigr)G_H}. \tag{1$'$}\]

\emph{Proof.} Integrate \(m_1,\dots,m_n\) in turn. Each step applies the
Brauer--Klimyk rule, which holds with the extended characters
\(\chi_\nu\) for every \(\nu\in P\) (its proof uses only the Weyl
invariance of the weight multiset), and orthogonality; this gives
\(k_t(g)E\prod\chi_{\lambda_j}(m_j)=\sum_\mu d_\mu\sum_{\boldsymbol\kappa}\prod m(\kappa_j) e^{-\frac t2[s_1C_2(\mu)+\sum_j(s_{j+1}-s_j)C_2(\mu+\kappa_1+\dots+\kappa_j)]}\chi_{\mu+\sum\kappa_j}(g)\),
\(s_{n+1}=1\). The unfolding of step (ii) holds because the multiset of
tuples \(\boldsymbol\kappa\) is invariant under the simultaneous Weyl
action. Completing the square,
\(s_1|v|^2+\sum_j(s_{j+1}-s_j)|v+\kappa_1+\dots+\kappa_j|^2=|v+A|^2+Q\),
where \(Q\) is the covariance of the standard Brownian bridge at the
times \(s_j\), paired with the \(\kappa_j\). Steps (iv)--(vi) are
unchanged with \((1-s)\kappa\) replaced by \(A\) and \(s\kappa\) by
\(\sum_js_j\kappa_j\). \(\square\)

The polynomial ratio for a single image is
\(\pi(\frac it Z-(1-s)\kappa)/\pi(\frac itZ)=\prod_{\alpha>0}\bigl(1+i(1-s)t\langle\kappa,\alpha\rangle/\langle Z,\alpha\rangle\bigr)\),
the curvature correction; for \(SU(2)\) at \(s=\frac12\) it is the
factor \(1+ikt/(2\theta)\) behind the
\(-\frac t{2\theta}k\sin\frac{k\theta}2\) term of the \(SU(2)\) note.

\hypertarget{reductions-corollary-2}{%
\subsection{3. Reductions (Corollary 2)}\label{reductions-corollary-2}}

\emph{\(U(1)\).} \(P=Q^\vee=\mathbb Z\), \(\pi\equiv1\), \(\rho=0\),
\(\langle n,H\rangle=nH\), \(X=\theta\):
\(E\,e^{inm}=e^{-ts(1-s)n^2/2}\sum_He^{in(s\theta+2\pi(1-s)H)}G_H/\sum_HG_H\),
the characteristic function of the mixture over the winding \(H\), with
weights \(\propto e^{-(\theta-2\pi H)^2/(2t)}\), of wrapped Gaussians of
variance \(s(1-s)t\) centred at \(s\theta+2\pi(1-s)H\). This is
Corollary 5\(_s\)(a) of the series/parallel note.

\emph{\(SU(2)\), \(s=\frac12\).} The weights of spin \(J\) are
\(\kappa=k\alpha\), \(k=-J,\dots,J\), with \(|\kappa|^2=k^2\),
\(\langle\kappa,X\rangle=k\theta\),
\(\langle\kappa,w\alpha^\vee\rangle=2kw\), and
\(\pi(v)=2\langle v,\alpha\rangle\). The phase is \(e^{2\pi ikw}\):
\((-1)^w\) for half-integer \(J\) and \(1\) for integer \(J\). With
\(\Xi_{(\sigma)},\Theta_{(\sigma)}\) as in the \(SU(2)\) note, the
numerator of (1) is
\(\sum_ke^{-tk^2/8}e^{ik\theta/2}\bigl[\frac it\Xi_{(\sigma)}-\frac k2\Theta_{(\sigma)}\bigr]\)
and the denominator is \(\frac it\Xi_+\). Symmetrizing in \(k\) gives
Theorem 3 of that note, and \(J=\frac12\) gives its Theorem 1,
\(e^{-t/32}[2\cos\frac\theta4\,\Xi_--\frac t2\sin\frac\theta4\,\Theta_-]/\Xi_+\).

\hypertarget{the-centre-corollary-3}{%
\subsection{4. The centre (Corollary 3)}\label{the-centre-corollary-3}}

\emph{Proof.} Two weights of one representation differ by an element of
the root lattice \(Q\), so \(e^{2\pi i(1-s)\langle\kappa,H\rangle}\) is
independent of \(\kappa\) for every \(\lambda\) iff
\(\langle\beta,(1-s)H\rangle\in\mathbb Z\) for all \(\beta\in Q\),
i.e.~iff \((1-s)H\in Q^*=P^\vee\). Then
\(e^{2\pi i(1-s)\langle\kappa,H\rangle}=e^{2\pi i\langle\lambda,(1-s)H\rangle}\)
is the scalar by which the central element \(z_H=e^{2\pi i(1-s)H}\) acts
in \(\lambda\). For \(s=p/q\) in lowest terms, \((1-s)=(q-p)/q\) with
\(\gcd(q-p,q)=1\), and \((q-p)H/q\in P^\vee\) iff \(H/q\in P^\vee\)
(Bézout, since \(H\in Q^\vee\subset P^\vee\)). So the central images are
\(H\in qP^\vee\cap Q^\vee\), and the elements reached are
\(e^{2\pi i(q-p)Y}\) with \(Y\in P^\vee\cap q^{-1}Q^\vee\): the
\(q\)-torsion of \(P^\vee/Q^\vee\cong Z(G)\), on which multiplication by
\(q-p\) is an automorphism. For \(SU(N)\), \(Z(G)=\mathbb Z_N\) and its
\(q\)-torsion is \(\mathbb Z_{\gcd(q,N)}\). \(\square\)

Consequences. Dyadic refinement (\(q=2^n\)) reaches the \(2\)-primary
part of the centre: all of it for \(SU(2)\) at the first halving,
\(\pm1\) at a halving and \(\mathbb Z_4\) at quarters for \(SU(4)\),
nothing for \(SU(3)\) or any odd \(N\). The image classes of a halving
are all central only for \(SU(2)\), where \(\frac12Q^\vee=P^\vee\).

Geometrically the central elements are focal points of \(e\): the
shortest geodesics from \(e\) to \(e^{2\pi iY}\) (\(Y\in P^\vee\) of
minimal norm in its class) form the adjoint orbit of \(Y\), a sphere
\(S^2\) for \(-1\in SU(2)\) and a \(\mathbb{CP}^2\) for
\(\omega\cdot1\in SU(3)\), because conjugation fixes both endpoints.

\hypertarget{su3-corollary-4}{%
\subsection{\texorpdfstring{5. \(SU(3)\) (Corollary
4)}{5. SU(3) (Corollary 4)}}\label{su3-corollary-4}}

Use \(\mathfrak t=\{H\in\mathbb R^3:\sum H_i=0\}\),
\(Q^\vee=\mathfrak t\cap\mathbb Z^3\), \(|H|^2=2\sum H_i^2\). The
weights of \(\mathbf3\) are \(\kappa_i=e_i-\frac13(1,1,1)\), so
\(\langle\kappa_i,H\rangle=H_i\): the image \(H\) carries the phase
\(e^{2\pi i(1-s)H_i}\) on the \(i\)-th weight.

\begin{itemize}
\tightlist
\item
  \(s=\frac13\): the phases \(e^{4\pi iH_i/3}\) agree iff
  \(H_1\equiv H_2\equiv H_3\pmod3\), i.e. \(H\in3P^\vee\). For
  \((1,1,-2)\) and its Weyl images all \(H_i\equiv1\) and the phase is
  \(e^{4\pi i/3}=\omega^2\); for \((-1,-1,2)\) and its images it is
  \(\omega\). The antifundamental gets the conjugates; the adjoint, with
  weights \(e_i-e_j\), gets \(e^{4\pi i(H_i-H_j)/3}=1\). These images
  have \(|H|^2=12\), so
  \(G_H/G_0=e^{-24\pi^2/t+2\pi\langle X,H\rangle/t}\). The root images,
  such as \((1,-1,0)\) with \(|H|^2=4\) and
  \(G_H/G_0\approx e^{-8\pi^2/t}\), give the fundamental the phases
  \(\omega^2,\omega,1\): weight-dependent.
\item
  \(s=\frac12\): the phases \(e^{\pi iH_i}\) agree iff all \(H_i\) have
  the same parity; with \(\sum H_i=0\) they are then all even,
  \(H\in2Q^\vee\), and the phase is \(1\).
\end{itemize}

So an \(SU(3)\) halving produces only weight-dependent phases
(non-central shifts), and a trisection produces the triality character
on the classes \(3P^\vee\), below the root images by the factor
\(e^{-16\pi^2/t}\).

\emph{Remark (a trisection as three fusing vortices; semiclassical
placement).} A trisection inserts two layers, at \(s=\frac13\) and
\(\frac23\); the marginal law of each is (1). For a central image
\(H=3Y\), \(Y\in P^\vee\), \(z=e^{2\pi iY}\), the image path
\(e^{iu(X-2\pi H)}\) passes the two layers at \(e^{iX/3}z^{-1}\) and
\(e^{2iX/3}z^{-2}\), so each of the three sub-segments carries the
increment \(e^{iX/3}z^{-1}\): every sub-face of the cut face gains the
same central flux \(z^{-1}\) relative to the direct path, and the three
fuse to \(z^{-3}=1\). The \(SU(2)\) halving is the two-piece case, where
each half-face gains \(-1\) and the two fuse to \(1\), Dirac's belt. For
\(SU(3)\) it is three thin \(\mathbb Z_3\) vortices in one cut face
fusing to the identity, the same bookkeeping by which three quarks form
a colour singlet (a parallel, not a derivation). Theorem 1\('\) makes
the placement exact at the level of character moments: for
\(s_1=\frac13\), \(s_2=\frac23\) and a central image \(H=3Y\) the phase
is
\(e^{2\pi i(2\langle\kappa_1,Y\rangle+\langle\kappa_2,Y\rangle)}=\zeta_{\lambda_1}(z)^2\zeta_{\lambda_2}(z)\)
with \(\zeta_\lambda(z)=e^{2\pi i\langle\lambda,Y\rangle}\) the central
character, which is the phase of the layers \(z^{-1}\) and \(z^{-2}\)
since \(\zeta^3=1\). Joint moments of class functions do not determine
the full joint law of the two layers, whose relative orientation remains
to be described.

\hypertarget{what-the-result-says-and-what-it-leaves-open}{%
\subsection{6. What the result says, and what it leaves
open}\label{what-the-result-says-and-what-it-leaves-open}}

It is an exact statement about one bridge, the building block of the
parallel insertion, for every compact simply connected group, every
representation and every cut fraction. It concerns characters only: the
diagonal entries of the midpoint matrix \(E\,D^\lambda(m)\) differ from
the weight terms of (1), since Theorem 4 of the \(SU(2)\) note gives the
spin-1 entries an additional Gaussian tail, so atlas cell 1 is
unaffected. In the refinement it fixes which windings of a single step
act through the centre. The weights are small per step: in \(1+3\) the
image exponents \(8\pi^2\) (roots) and \(24\pi^2\) (central \(SU(3)\)
trisection images) lie far beyond the per-volume threshold
\(c>2/b_0=32\pi^2/11\) of the
\href{zero-spacing-any-action.md}{zero-spacing note} for \(SU(3)\).
Centre-valued link configurations exist on any lattice whatever the
blocking, and thick centre vortices
(\href{https://doi.org/10.1016/0003-4916(79)90346-4}{Mack and Petkova
1979}; \href{https://doi.org/10.1016/0550-3213(78)90153-0}{'t Hooft
1978}; metadata) are a question about many steps. The centre matters
because pure \(SU(3)\) leaves it unscreened: by the complementarity of
\href{https://doi.org/10.1103/PhysRevD.19.3682}{Fradkin and Shenker
(1979)} and
\href{https://doi.org/10.1016/0003-4916(78)90039-8}{Osterwalder and
Seiler (1978)} (metadata), matter in the fundamental representation
screens every centre charge and joins the confinement and Higgs regimes
without a transition, while in the pure gauge theory the unbroken
\(\mathbb Z_3\) stays an order parameter. The centre-reaching (triadic)
windings are therefore a feature of the pure-gauge goal, screened in QCD
with quarks, which is one more difference between the continuum limits
compared in the \href{three-continuum-limits.md}{joint-paper plan}.
Whether a triadic refinement, \(b=3\), organizes the \(SU(3)\)
large-field terms better than dyadic refinement is open; factor-2
decimation with the \(\mathbb Z_2\) factor kept explicit goes back to
\href{https://doi.org/10.1103/PhysRevD.23.2371}{Tomboulis (1981)}
(metadata), and prior art on \(b=3\) or centre-adapted blocking for
\(SU(3)\) has not been searched.

\hypertarget{consequence-for-state}{%
\subsection{7. Consequence for STATE}\label{consequence-for-state}}

Atlas §1b: the bullet on which part of the centre a cut reaches now
rests on Corollary 3 of an exact formula, (1), which also carries the
\(SU(2)\) belt trick and the \(SU(3)\) triality at thirds. Refereeing by
Astra is pending; the exact trisection formula for the fundamental of
\(SU(3)\) is (1) with \(s=\frac13\) and the weights of §5.
\end{document}
